\documentclass[10pt]{article}
\usepackage{icmlko}

\icmlmeta{IP 파운데이션 월드모델}{328}{우리가 안 잰 바닥}{쉰 노트를 유보 재채점 자로 쟀는데, 학습 쪽 흔들림이 두 배에서 여덟 배였다}

\begin{document}

\twocolumn[
\icmlheader

\begin{abstract}
노트 274가 검출 바닥을 $2{\times}$짝SE 로 놓았고 그 뒤 쉰 노트가 그 자로
쟀다. 짝SE 는 \textbf{유보를 다시 뽑아 다시 채점하면} 얼마나 흔들리나다.
노트 327이 아이돌에서 \emph{학습 쪽}을 흔들어 보니 1.65배였다.
\textbf{오늘 열 도메인 전부 쟀다} --- 학습 행을 재추출해 다시 적합하고
고정된 유보를 채점, 뽑기 서른.

\textbf{학습SE 가 짝SE 의 2배에서 8배다.} 그리고 \textbf{비가 큰 도메인일수록
크다}: 웹툰 $7.99{\times}$ · 모바일 $7.48{\times}$ · 애니 $6.53{\times}$ ·
세계애니 $4.01{\times}$ · 도서 $2.47{\times}$ · 아이돌 $2.13{\times}$.
\textbf{채점 잡음은 표본 수를 따르는데 학습 잡음은 안 따른다} --- 짝SE 는
0.004$\sim$0.056 으로 14배 벌어져 있고 학습SE 는 0.025$\sim$0.119 로
5배뿐이다. 큰 도메인의 ``0.004'' 는 정밀한 게 아니라 \textbf{다른 것을 잰
것}이었다.

결과가 셋이다. ① \textbf{판 자체가 학습 뽑기마다 $\pm$0.031 움직인다} ---
배포 규칙의 미결정 폭 0.010 의 \textbf{3.1배}. ② \textbf{도메인 순서가
안 선다.} 같은 여섯으로 견주면 짝SE 로는 15쌍 중 14쌍(93\%)이 서고 이웃
4/5 인데, 학습SE 로는 10쌍(67\%)이고 \textbf{이웃 0/5} --- 붙어 있는 두
도메인의 앞뒤를 하나도 주장할 수 없다. ③ \textbf{그런데 짝지으면 갈린다.} 같은 뽑기에서 두 팔을 다
돌리면 --- \emph{축만} 바꾼 짝(\texttt{fund\_cat})은 SD 0.0203 에 부호가
\textbf{14/14} 로 안 뒤집히고, \emph{행까지} 바꾼 짝(노트 327의 아이돌
넓힘)은 SD 0.1308 에 \textbf{20 중 3이 음수}다. \textbf{지난 쉰 노트는
거의 다 축만 바꿨으므로 방향은 맞았을 것이다} --- 다만 미결정 폭 0.010 은
짝지은 SD 0.0203 의 절반이라 \textbf{``갈렸다''를 너무 자주 말했다.}
\end{abstract}
]

\section{자가 셋인데 하나만 써 왔다}

\begin{center}\footnotesize
\begin{tabular}{lll}
\hline
자 & 무엇을 흔드나 & 아이돌 SD \\
\hline
짝SE (노트 274 · 316) & 유보 25행 재채점 & 0.056 \\
위약 (노트 327) & 학습 라벨 섞기 & 0.092 \\
\textbf{학습 재추출} (오늘) & \textbf{학습 행 재추출} & \textbf{0.119} \\
\hline
\end{tabular}
\end{center}

\begin{figure*}[t]
\centering
\includegraphics[width=\textwidth]{figs/floor.pdf}
\caption{왼쪽 --- 도메인별 두 SE. 가운데 --- 학습 뽑기 서른의 판.
오른쪽 --- 순서가 몇 쌍이나 서나.}
\end{figure*}

\section{열 도메인}

한 뽑기마다 \textbf{모든 도메인의 학습 행}을 재추출하고 다시 적합해 고정된
유보를 채점한다. F18 배깅부스트 · \texttt{trendcalwikitag} · $t15$ · 서른 번.

\begin{center}\footnotesize
\begin{tabular}{lrrrr}
\hline
도메인 & 기준 & 학습SE & 짝SE & 비 \\
\hline
게임 & 0.5763 & 0.0538 & --- & --- \\
모바일 & 0.5439 & 0.0299 & 0.004 & \textbf{7.48} \\
세계애니 & 0.5408 & 0.0441 & 0.011 & 4.01 \\
애니 & 0.4701 & 0.0392 & 0.006 & \textbf{6.53} \\
웹툰 & 0.4480 & 0.0320 & 0.004 & \textbf{7.99} \\
팝업 & 0.4346 & 0.0447 & --- & --- \\
도서 & 0.3616 & 0.0740 & 0.030 & 2.47 \\
시장팝업 & 0.2869 & 0.0635 & --- & --- \\
펀딩 & 0.2556 & 0.0245 & --- & --- \\
아이돌 & 0.2108 & \textbf{0.1192} & 0.056 & 2.13 \\
\hline
\textbf{판} & \textbf{0.3030} & \textbf{0.0308} & & \\
\hline
\end{tabular}
\end{center}

\textbf{비가 큰 도메인일수록 크다.} 짝SE 는 $1/\sqrt{n}$ 을 따라 웹툰
0.004 까지 내려가는데 학습SE 는 안 내려간다 --- 학습이 2,106행이어도
0.032 다. \textbf{어느 2,106행을 뽑았느냐가 유보 711행을 어떻게 채점하느냐
보다 훨씬 크게 결과를 흔든다.}

\section{판 미결정 폭 0.010}

노트 275 · 282가 배포 규칙을 세웠다 --- ``판이 0.010 미결정 폭 안이면
배포가 정한다''. 그 0.010 은 짝SE 계열의 수다.

\begin{center}\small
\begin{tabular}{lr}
\hline
판 학습SE (뽑기 서른) & \textbf{0.0308} \\
판 값의 범위 & [0.2349, 0.3778] \\
미결정 폭 & 0.010 \\
\hline
비 & \textbf{3.1배} \\
\hline
\end{tabular}
\end{center}

\section{순서가 안 선다}

\texttt{lab/ordering.py}(노트 315)로 도메인 순서에 몇 쌍이 서는지 셌다.
짝SE 가 있는 여섯으로 견준다.

\begin{center}\footnotesize
\begin{tabular}{lrr}
\hline
쓴 SE & 선 쌍 & 이웃 쌍 \\
\hline
짝SE (노트 316) & 14/15 (93\%) & \textbf{4/5} \\
학습SE & 10/15 (67\%) & \textbf{0/5} \\
둘을 합침 & 9/15 (60\%) & \textbf{0/5} \\
\hline
\end{tabular}
\end{center}

\textbf{이웃 쌍이 하나도 안 선다.} 열 도메인으로 넓혀도 45쌍 중 29(64\%)에
이웃 0/9 다. 노트 314가 ``못 가르는 값들의 순서도 못 가른다''를 배웠고
노트 316이 ``짝SE 가 14배 벌어져 도메인 순서를 못 읽는다''로 이어졌는데,
\textbf{그때 쓴 자가 셋 중 제일 작은 것이었다.}

\section{짝지어 재도 그렇다}

흔들림 대부분이 뽑기마다 \emph{공통}이면 두 팔의 \textbf{차}는 안정할 수
있다. 그래서 노트 327의 아이돌 넓힘을 \textbf{같은 뽑기에서 두 팔 다}
돌렸다(뽑기 스물).

\begin{center}\footnotesize
\begin{tabular}{lrrr}
\hline
 & 평균 차 & SD & 양수 \\
\hline
아이돌 & $+$0.1527 & \textbf{0.1308} & \textbf{17/20} \\
판 & $+$0.0126 & 0.0285 & --- \\
\hline
\multicolumn{4}{l}{재추출 없이 잰 것: 아이돌 $+$0.1862 · 판 $+$0.0031} \\
\hline
\end{tabular}
\end{center}

\textbf{평균은 확실하다}($t{=}5.22$). \textbf{한 번의 관측은 안 확실하다}
--- SD 0.1308 이고 스물 중 셋이 음수다. 우리가 실제로 보는 것은 뽑기
\emph{하나}다. 그리고 판의 짝지은 차 SD 0.0285 도 미결정 폭 0.010 의
\textbf{2.9배}다.

\section{그런데 축만 바꾸면 안 흔들린다}

지금까지 판을 견뎌 온 방식은 \textbf{행이 같고 축만 다른} 짝이다. 그것도
같은 뽑기에서 둘 다 돌렸다(\texttt{trendcalwikitag} 대 $+$\texttt{fund\_cat},
뽑기 열넷).

\begin{center}\footnotesize
\begin{tabular}{lrrr}
\hline
바꾼 것 & 평균 차 & SD & 양수 \\
\hline
\textbf{축만}(판) & $+$0.0294 & \textbf{0.0203} & \textbf{14/14} \\
\quad 펀딩 & $+$0.0819 & 0.0358 & 14/14 \\
\hline
\textbf{행까지}(아이돌) & $+$0.1527 & \textbf{0.1308} & 17/20 \\
\quad 판 & $+$0.0126 & 0.0285 & --- \\
\hline
\end{tabular}
\end{center}

\textbf{부호가 열넷에 열넷이다.} 축을 더하는 처치는 뽑기가 바뀌어도 방향이
안 뒤집힌다 --- 두 팔이 \emph{같은 행}을 보므로 흔들림이 공통이고 차에서
지워진다. \textbf{행을 바꾸는 처치는 그게 안 된다.}

\begin{center}\small
\begin{tabular}{ll}
\hline
축만 바꾼다 & 부호는 믿어라 · \textbf{크기는 $\pm$0.02} \\
행을 바꾼다 & \textbf{부호도 다시 봐라} ($\pm$0.13) \\
\hline
\end{tabular}
\end{center}

\textbf{그래서 지난 쉰 노트의 결론 대부분은 방향이 맞았을 것이다} --- 거의
다 축만 바꾼 짝이었다. 다만 \textbf{미결정 폭 0.010 은 여전히 절반이다.}
$\Delta{=}0.015$ 를 ``갈렸다''고 부른 자리는 사실 $0.75\sigma$ 였다.
\textbf{폭을 0.02 로 넓히는 것이 오늘 잰 것에 맞는다.}

\paragraph{덤 --- 안 맞는 수 하나} \texttt{fund\_cat} 이 여기서 판
$+$0.0289 인데 노트 309는 $+$0.0023 으로 적었다. 그 사이에 시장팝업이
채점 도메인으로 들어왔고(노트 284) 위키 표시자가 정정됐다(노트 306).
\textbf{어느 쪽이든 옛 수와 새 수를 직접 견주면 안 된다는 규약이 또
필요했던 자리다.} 다시 재야 한다.


\section{남은 것}
\begin{center}\footnotesize
\begin{tabular}{ll}
\hline
갈래 & 왜 \\
\hline
\textbf{지난 판정을 다시 걸러} & 바닥 아래였던 것이 몇인가 \\
미결정 폭을 다시 정할까 & 0.010 $\to$ 0.03 이면 거의 다 미결정 \\
학습SE 를 매 실행에 & 청력 · 겹침 같은 이름표로 \\
왜 큰 도메인이 안 줄까 & 축이 아니라 계수가 흔들리나 \\
\hline
\end{tabular}
\end{center}

\paragraph{규약} \textbf{자는 잰 것만 잰다.} 짝SE 는 ``같은 모형으로 다른
유보를 채점하면''을 묻고, 그 답이 작다고 해서 ``이 수가 믿을 만하다''가
되지 않는다. 우리가 하려는 말은 늘 \textbf{다른 자료에서도 이럴까}인데
그건 학습 쪽을 흔들어야 나오는 답이다. \textbf{쉰 노트 동안 한쪽만 흔들고
있었다.} 그리고 이 노트가 뒤집는 것은 결론이 아니라 \emph{확신의 크기}다
--- 오늘까지의 판정 대부분은 방향이 맞았을 것이다. 다만
\textbf{``$t{=}4.6$ 이라 확실하다''고 쓴 자리들은 다시 봐야 한다.}

\end{document}
